% Generated by tools/textbook/export_listings.js from tools/concepts/assessments.json: do not edit.
\begin{chapstatus}{Il filo del capitolo}
Un programma è fatto di poche funzioni che il runtime chiama; ogni riga è un'istruzione; i valori stanno nei registri mentre si calcola e in memoria tra un frame e l'altro. Tutto il resto che avete usato oggi ha una sua lezione. Prossimo passo: da dove vengono queste idee.
\par\medskip\textbf{Biglietto d'uscita} (da due a quattro minuti, da soli, su carta):
\begin{enumerate}[nosep,leftmargin=2em]
\item Leggete ad alta voce la riga  loop:  li a0, 32   \# x . Nominatene etichetta, istruzione, operandi e commento. li è un'istruzione del processore? Che cosa ne fa l'assemblatore?
\item Il debugger mostra l'istruzione addi t1, t1, 1 conservata in memoria come la parola 0x00130313. Che cos'è quella parola, e che cosa vi dice su dove abita un programma?
\end{enumerate}
\end{chapstatus}
